\documentclass[10pt,letterpaper,twocolumn]{article}
\usepackage[left=0.875in,right=0.875in,top=0.82in,bottom=0.75in,headheight=14pt,headsep=12pt]{geometry}
\usepackage{amsmath,amssymb,graphicx,booktabs,caption,fancyhdr,array,enumitem}
\usepackage[colorlinks=true,urlcolor=blue]{hyperref}
\hypersetup{pdftitle={JOREK ES to EM OpsSplit: Results and Methods},pdfauthor={JOREK OpsSplit study}}
\setlength{\columnsep}{0.25in}
\setlength{\parindent}{0pt}
\setlength{\parskip}{5pt}
\pagestyle{fancy}\fancyhf{}
\fancyhead[C]{\small\bfseries Physical Operator Learning in Gridded JOREK}
\fancyfoot[L]{\scriptsize ES--EM transfer\enspace---\enspace results 4 November 2026; methods expanded 9 November 2026; evidence snapshot \SnapshotTime}
\fancyfoot[R]{\scriptsize\thepage/2}
\renewcommand{\headrulewidth}{0.4pt}
\captionsetup{font=small,labelfont=normalfont}
\makeatletter
\renewcommand\section{\@startsection{section}{1}{0pt}{8pt}{4pt}{\normalfont\fontsize{11}{13}\bfseries}}
\makeatother
\input{experiment_values.tex}
\newcount\companionpages
\AddToHook{shipout/before}{\global\advance\companionpages by 1
\ifnum\companionpages>2\PackageError{pagelimit}{TWO PAGE LIMIT EXCEEDED}{Shorten the content.}\fi}
\begin{document}
\twocolumn[{
\begin{center}
{\large\bfseries JOREK OpsSplit: Electrostatic-to-Electromagnetic Transfer}\par
\vspace{3pt}{\small Model003 $\rightarrow$ model600\enspace---\enspace full-data, three-seed experiment}
\end{center}\vspace{-3pt}
\textbf{Result.} In this configuration, ES pretraining gives no consistent advantage over training the same split from scratch.
Long rollouts expose large errors despite finite outputs. This tests an effective learned split, not recovery of individual physical terms.
\begin{center}\includegraphics[width=\textwidth]{results_figure.pdf}\end{center}\vspace{-10pt}
\captionof{figure}{Ten shared EM holdouts; lower is better. Dashed: persistence. Panel (b): seed means; zero-shot uses no EM training.}\vspace{5pt}
}]
\section*{1\quad Data and controlled protocol}
The releases contain 497 ES and 100 EM trajectories on a $100\!\times\!100$ sampled $(R,Z)$ grid,
with 201 (ES) or 211 (EM) saved frames each.
Train/validation/test: ES 400/44/53; EM 76/14/10.
Splits are assigned by a hash of each run's blob geometry, so blob-family hashes prevent cross-regime split leakage.
Time is the saved-frame index.

Each seed (0, 1, 2) retrains the ES source and every EM arm.
EM training subsets are nested: cases are ordered by a hash of (seed, family, id) and the first 8/32/64 are taken,
so $8\subset32\subset64$ within a seed, while different seeds draw different subsets.
Within a seed, all arms share these subsets, the 14 validation and the 10 test trajectories.

% Defined on page 1 so the full-width table lands at the top of page 2.
\begin{table*}[t]
\centering\small
\caption*{Table 1: How each EM arm is conducted. ``ES'': copied from that seed's ES source checkpoint; ``EM data'': EM training trajectories.
Parameter counts, learning rates and times come from run provenance. Train: wall time of the 20 epochs with validation on one H100 GPU,
mean over three seeds and budgets (budget-independent), excluding test rollouts; ``$+$\SourceMinutes'': the reused ES source pretraining.}
\vspace{-4pt}
\setlength{\tabcolsep}{4pt}
\begin{tabular}{@{}lllrcccr@{}}\toprule
Arm & Fluid $F$ & Correction $M_\varphi$ & Trainable / total & LR & Scaling & EM data & Train (min)\\\midrule
Zero-shot & ES, frozen & removed; $\psi,j$ persist & \TrainZeroShot\ / \TotalZeroShot & -- & ES & none & 0 $+$\SourceMinutes\\
Frozen transfer & ES, frozen & random, trained & \TrainFrozen\ / \TotalFrozen & $\LrFrozen$ & ES & 8/32/64 & \MinutesFrozen\ $+$\SourceMinutes\\
Fine-tuned & ES init, trained & random, trained & \TrainFinetune\ / \TotalFinetune & $\LrFinetune$ & ES & 8/32/64 & \MinutesFinetune\ $+$\SourceMinutes\\
Scratch split & random, trained & random, trained & \TrainScratch\ / \TotalScratch & $\LrScratch$ & EM & 8/32/64 & \MinutesScratch\\
Monolithic & \multicolumn{2}{l}{single $W_\varphi(q)$, random, trained} & \TrainMonolithic\ / \TotalMonolithic & $\LrMonolithic$ & EM & 8/32/64 & \MinutesMonolithic\\
\bottomrule\end{tabular}
\vspace{-6pt}
\end{table*}

\section*{2\quad Additive operator formulation}
For standardized fields $q_f=(\rho,T,u,\omega)$ and $q_b=(\psi,j)$, learn
\begin{equation*}
\begin{aligned}
F_\theta(q_f)&=N_\theta(q_f)+\operatorname{diag}(d_\theta)\Delta_{R,Z}q_f,\\
G_{\theta,\varphi}(q)&=(F_\theta(q_f),0)+M_\varphi(q),\\
q^{n+1}&=q^n+G_{\theta,\varphi}(q^n),\\
\Delta_{R,Z}&=\partial_{RR}+R^{-1}\partial_R+\partial_{ZZ}.
\end{aligned}
\end{equation*}
$q_f$ holds the four fields both regimes share; $q_b$ holds the EM-only flux and current.
$N$ is the fluid FNO ($4\!\to\!4$), and it sees only $q_f$.
FNO: width 16, 12 modes/axis, three layers, 20\% domain padding.
The diffusivity $d_\theta=d_{\max}\,\sigma(\ell_\theta)$ is learned per field.
Its cap $d_{\max}=0.2/(\Delta R^{-2}+\Delta Z^{-2})$ keeps this positive effective diffusion below the explicit-Euler limit, and $\Delta_{R,Z}$ acts on interior points only; neural stability is not guaranteed.
$M$ is the EM correction, a $6\!\to\!6$ FNO of the same size.
It sees all six fields and writes all six channels, so it supplies the magnetic tendency and can also correct the fluid one.
One step advances one saved frame.

\section*{3\quad How each OpsSplit arm is conducted}
\textbf{Stage 1: ES source.} For each seed, a four-channel split $F_\theta$ (no $M$) is trained on all 400 ES training
trajectories with the loop of Section~4. On ES data the third channel is the electrostatic potential $\Phi$.
The checkpoint $\theta_{\rm ES}$ ($N_\theta$ and $d_\theta$, \TrainSource\ parameters) is kept along with the ES per-field mean and SD.

\textbf{Stage 2: five EM arms.} Table~1 lists what each arm loads, trains and rescales.
The arms are built as paired contrasts.
\begin{itemize}[leftmargin=10pt,itemsep=1pt,topsep=2pt,parsep=0pt]
\item \emph{Zero-shot vs.\ persistence:} does the ES fluid operator alone predict EM fluid evolution?
\item \emph{Frozen vs.\ fine-tuned:} can the ES fluid operator be reused as is, or must it adapt?
\item \emph{Fine-tuned vs.\ scratch split:} same architecture, capacity, data and update count, but different initialization.
The learning rate and shared-field scaling also differ, so this is not a single-variable contrast.
\item \emph{Scratch split vs.\ monolithic:} does the split structure help compared with one unstructured operator?
\end{itemize}

\textbf{Zero-shot.} $\theta_{\rm ES}$ is copied into $F$ and nothing is trained.
The step runs with $M$ removed: $q_f$ advances under $F_{\theta_{\rm ES}}$, while $\psi$ and $j$ keep their frame-0 values.
No EM trajectory or EM statistic is used.

\textbf{Frozen transfer.} $\theta_{\rm ES}$ is copied and frozen; only a randomly initialized $M_\varphi$ trains.
$M$ must absorb induction, fluid--magnetic coupling and any ES-to-EM fluid mismatch.
$d_\theta$ stays at its ES values.

\textbf{Fine-tuned.} Same initialization, but $N_\theta$, $d_\theta$ and $M_\varphi$ all train, at a lower learning rate so the fluid part stays close to $\theta_{\rm ES}$.

\textbf{Scratch split.} The fine-tuned architecture trained from random initialization, with EM statistics for all six fields.

\textbf{Monolithic.} A single FNO $W_\varphi$ ($6\!\to\!6$) with $q^{n+1}=q^n+W_\varphi(q^n)$: no split and no diffusion term.

\textbf{Scaling in transfer.} Transfer arms keep the ES mean/SD for $(\rho,T,u,\omega)$, so the copied network receives the shared fields in the units it was trained on.
Magnetic scaling uses EM training data only.
The third channel, however, is $\Phi$ in ES and $u$ in EM.
Under ES scaling, EM $u$ has a standardized SD of 0.17--0.24, and $\omega$ has 0.4--0.6, rather than 1 (every seed and budget).
The copied operator therefore sees these two fields at a fraction of the amplitude it was trained on.

\section*{4\quad Common training, selection and evaluation}
\textbf{Training.} Every trained arm, and the ES source, uses the same loop.
A sample is a five-frame window $(q^s,\dots,q^{s+\Horizon})$ drawn uniformly, with replacement, from all start frames of the training trajectories.
The model rolls out freely from $q^s$, and the loss is the standardized squared error averaged over the \Horizon\ steps:
\begin{equation*}
\mathcal L=\frac{1}{\Horizon}\sum_{k=1}^{\Horizon}\big\langle\big(G^{k}(q^s)-q^{s+k}\big)^2\big\rangle,
\end{equation*}
where $\langle\cdot\rangle$ averages over batch, channels and grid.
For EM arms, $\psi$ and $j$ enter $\mathcal L$ with the same weight as the shared fields.
Optimizer: AdamW (weight decay $10^{-4}$), gradient-norm clip 1, batch eight.
Training runs \Epochs\ epochs of \StepsPerEpoch\ steps, i.e.\ \Updates\ updates whatever the budget.
With 8 EM cases ($\approx$1.7k windows), each window is drawn about 12 times on average; with 64 cases ($\approx$13k windows), about 1.5 times.

\textbf{Selection.} After each epoch, the same loss on the validation trajectories (windows every 20 frames) picks the best checkpoint.
EM arms always validate on all 14 EM validation trajectories.
At the 8-case budget, model selection therefore sees more EM trajectories than training does.

\textbf{Testing.} From frame 0 of each of the 10 EM holdouts, each arm rolls out 100 saved frames autoregressively, without resetting to the truth.
For field $f$ at horizon $h$, with prediction $\tilde q$ and sums over the grid,
\begin{equation*}
E_{h,f}=\Big(\textstyle\sum R\,(\tilde q^{\,h}_f-q^{h}_f)^2\big/\sum R\,(q^{h}_f)^2\Big)^{1/2}
\end{equation*}
in unstandardized units; $R$ is the toroidal-volume weight.
$E_h$ averages $E_{h,f}$ over $(\rho,T,u,\omega)$ and the ten holdouts.
$\psi$ and $j$ are trained but not scored; persistence predicts $q^0$ at every horizon.

\section*{5\quad Full-study results}
\begin{center}\small\setlength{\tabcolsep}{2pt}
\begin{tabular}{lccc}\toprule
EM cases & 8 & 32 & 64\\\midrule
\ResultsRows
\bottomrule\end{tabular}\\[2pt]
{\small Table 2: $E_{10}$, mean $\pm$ sample SD across three seeds.}\end{center}\vspace{-4pt}
Zero-shot: $E_{10}=\ZeroTen$ and $E_{100}=\ZeroHundred$; persistence: \PersistenceTen\ and \PersistenceHundred.
All rollouts remain finite, yet the worst eight-case fine-tuning seed reaches $E_{100}=\texttt{\WorstEight}$.
Seeds share ten test cases, not 30 independent holdouts.

\section*{6\quad Interpretation and next checks}
ES has no induction targets: full-EM zero-shot is not established.
The correction absorbs magnetic effects and mismatch; endpoint data do not identify PDE terms.
Potential/vorticity and flux/current constraints are not enforced.

Export units and timestamps remain unaudited.
Next: audit exports/constraints, test transfer with regime-appropriate scaling of $u$ and $\omega$, then select longer-rollout training on validation data.

{\footnotesize\textbf{Sources.} \href{https://zenodo.org/records/11099659}{ES: Zenodo 11099659};
\href{https://zenodo.org/records/11099527}{EM: 11099527};
\href{https://doi.org/10.1016/j.cpc.2024.109391}{Pamela et al., CPC 307, 109391}.}
\end{document}
